% !TEX spellcheck = en_US



% ~~~~~~~~~~~~~~~~~~~~~~~~~~~~~~~~~~~~~~~ V
% (NC) 2026 Haluk Bingol
% github.com/halukbingol/LaTeX-Templates
%
% Licensees may copy, distribute, display, and perform the work and make derivative works 
% and remixes based on it only for non-commercial purposes. 
% ~~~~~~~~~~~~~~~~~~~~~~~~~~~~~~~~~~~~~~~ A




% =====================================================================
%  Linked Lists -- Senior Level
%  ---------------------------------------------------------------------
%  Built on hblib/Template/hblibTemplate.tex with the libraries in ../../hblib/.
%  Programs and their real outputs are in codeoutput/:
%     codeoutput/<Name>.java   a program (classes without main() are only compiled)
%     codeoutput/<Name>.txt    what it printed
%     codeoutput/Javap*.txt    captured with javap -p on the JDK in this build
% =====================================================================

\documentclass[aspectratio=169,10pt,compress]{beamer}

% ~~~~~~~~~~~~~~~~~~~~~~~~~~~~~~~~~~~~~~~~~~~~~~~~~~~~~~~~~~~~~~~~~ V hblib
	\usepackage{../../hblib/hblibPresentation} % lib presentation: theme, i/N, \hSection, algorithm2e, ...
	\usepackage{../../hblib/hblibCodeoutput} % lib code + output: \lstcode, \lstterminal, ...
	\usepackage{../../hblib/hblibDatastructures} % lib data structures: \hString, \hStack, \hQueue
% ~~~~~~~~~~~~~~~~~~~~~~~~~~~~~~~~~~~~~~~~~~~~~~~~~~~~~~~~~~~~~~~~~ A hblib


\usepackage[T1]{fontenc}
\usepackage{lmodern}
\usepackage{booktabs}

\newcommand{\win}{\textbf{Windows}}
\newcommand{\mac}{\textbf{macOS}}

% long javap lines: wrap instead of running off the frame
\lstset{columns=fixed, breaklines=true, breakatwhitespace=false}
\lstdefinelanguage{text}{}      % plain text: javap output

\SetAlCapFnt{\small}
\SetAlCapNameFnt{\small}
\SetKwProg{Fn}{function}{}{end}
\SetKw{KwError}{error}

% ---------------------------------------------------------------------
\title{Linked Lists}
\subtitle{Senior Level: JDK Internals, Sentinel Nodes, Memory}
\author[Bingol]{Ruhan Alpaydin}

\institute[FCIS]{
    Faculty of Computer and Information Sciences\\
    Yeditepe University
}
\date{
	{\tiny \hbTimeStamp}
}

\begin{document}




% ~~~~~~~~~~~~~~~~~~~~~~~~~~~~~~~~~~~~~~ V frm
\begin{frame}[t,fragile]
  \titlepage
\end{frame}
% ~~~~~~~~~~~~~~~~~~~~~~~~~~~~~~~~~~~~~~ A frm




% ~~~~~~~~~~~~~~~~~~~~~~~~~~~~~~~~~~~~~~ V frm
\begin{frame}[t,fragile] \frametitle{Outline}
  \begin{columns}[T]
    \begin{column}{0.48\textwidth}
      \tableofcontents[sections={1-3}, hideallsubsections]
    \end{column}
    \begin{column}{0.48\textwidth}
      \tableofcontents[sections={4-5}, hideallsubsections]
    \end{column}
  \end{columns}
\end{frame}
% ~~~~~~~~~~~~~~~~~~~~~~~~~~~~~~~~~~~~~~ A frm

% ~~~~~~~~~~~~~~~~~~~~~~~~~~~~~~~~~~~~~~ V
\hSection[Internals]{JDK Internals}
% ~~~~~~~~~~~~~~~~~~~~~~~~~~~~~~~~~~~~~~ A




% ~~~~~~~~~~~~~~~~~~~~~~~~~~~~~~~~~~~~~~ V frm
\begin{frame}[t,fragile] \frametitle{\texttt{javap}: \texttt{LinkedList}}
  \lstcode[basicstyle=\ttfamily\scriptsize, firstline=3, lastline=11, firstnumber=3]{codeoutput/JavapLinkedList.txt}{text}
  \begin{itemize}\small
    \item Three private fields: \hCode{size}, \hCode{first}, \hCode{last}.
    \item \hCode{first} and \hCode{last} are references to \hCode{Node<E>},
          not an array. Nothing else is stored.
    \item \hCode{size} is cached: \hCode{size()} is $O(1)$.
    \item \hRemark{The class shape tells you the whole design}:
          a doubly linked list with a cached length.
  \end{itemize}
\end{frame}
% ~~~~~~~~~~~~~~~~~~~~~~~~~~~~~~~~~~~~~~ A frm




% ~~~~~~~~~~~~~~~~~~~~~~~~~~~~~~~~~~~~~~ V frm
\begin{frame}[t,fragile] \frametitle{\texttt{javap}: \texttt{ArrayList} (for comparison)}
  \lstcode[basicstyle=\ttfamily\scriptsize, firstline=3, lastline=11, firstnumber=3]{codeoutput/JavapArrayList.txt}{text}
  \begin{itemize}\small
    \item \hCode{elementData} is a plain \hCode{Object[]}: the storage is one block.
    \item \hCode{size} is a counter, not an index.
    \item No \hCode{first} or \hCode{last}: there is one array, and access
          is \hCode{elementData[i]}.
    \item \hRemark{The two classes have nothing in common except the \texttt{List} interface}.
  \end{itemize}
\end{frame}
% ~~~~~~~~~~~~~~~~~~~~~~~~~~~~~~~~~~~~~~ A frm




% ~~~~~~~~~~~~~~~~~~~~~~~~~~~~~~~~~~~~~~ V frm
\begin{frame}[t,fragile] \frametitle{The \texttt{Node} class}
\begin{lstlisting}[language=Java, basicstyle=\ttfamily\small, frame=single, xleftmargin=2em]
private static class Node<E> {
    E item;
    Node<E> next;
    Node<E> prev;

    Node(Node<E> prev, E element, Node<E> next) {
        this.item = element;
        this.next = next;
        this.prev = prev;
    }
}
\end{lstlisting}
  \begin{itemize}\small
    \item \hCode{private}: the outside world never sees a \hCode{Node}.
    \item \hCode{static}: the node does not belong to one list;
          it is a standalone container that any \hCode{LinkedList} uses.
    \item Three fields plus one constructor. That is the whole class.
    \item \hRemark{The constructor sets all three at once}, so no
          half-built node ever exists.
  \end{itemize}
\end{frame}
% ~~~~~~~~~~~~~~~~~~~~~~~~~~~~~~~~~~~~~~ A frm




% ~~~~~~~~~~~~~~~~~~~~~~~~~~~~~~~~~~~~~~ V frm
\begin{frame}[t,fragile] \frametitle{Memory layout}
  \begin{columns}[T]
    \begin{column}{0.56\textwidth}
      \lstcode[basicstyle=\ttfamily\tiny]{codeoutput/MemoryLayout.java}{Java}
    \end{column}
    \begin{column}{0.40\textwidth}
      \lstterminal[basicstyle=\ttfamily\scriptsize]{codeoutput/MemoryLayout.txt}
      \small
      The measurement is rough: \hCode{Runtime} reports bytes with
      and without the boxed array. \hRemark{Roughly 20 bytes per boxed element},
      versus 4 bytes for \hCode{int[]}.
    \end{column}
  \end{columns}
  \vspace{0.3em}
  \small
  A \hCode{LinkedList<Integer>} adds \emph{two} references per node on
  top of that: each node is about $16 + 8 + 8$ bytes.
\end{frame}
% ~~~~~~~~~~~~~~~~~~~~~~~~~~~~~~~~~~~~~~ A frm

% ~~~~~~~~~~~~~~~~~~~~~~~~~~~~~~~~~~~~~~ V
\hSection[Sentinel]{Sentinel Nodes}
% ~~~~~~~~~~~~~~~~~~~~~~~~~~~~~~~~~~~~~~ A




% ~~~~~~~~~~~~~~~~~~~~~~~~~~~~~~~~~~~~~~ V frm
\begin{frame}[t,fragile] \frametitle{What a sentinel is}
  \begin{block}{The problem}
    \small
    Without a sentinel, every \hCode{addFirst} and \hCode{removeFirst}
    needs a check: \emph{is the list empty?} The first node is special
    because there is no ``previous node'' to relink.
  \end{block}
  \begin{block}{The sentinel fix}
    \small
    Keep a dummy node that never holds data. \hCode{head} always points
    at it. \hCode{head.next} is the first real node (or \hCode{null}).
    Now \hCode{addFirst} and \hCode{removeFirst} treat the empty list
    exactly like the non-empty one.
  \end{block}
  \vspace{0.2em}
  \small
  \hRemark{One extra node removes every special case}.
  The cost is one node of memory; the benefit is code without branches.
\end{frame}
% ~~~~~~~~~~~~~~~~~~~~~~~~~~~~~~~~~~~~~~ A frm




% ~~~~~~~~~~~~~~~~~~~~~~~~~~~~~~~~~~~~~~ V frm
\begin{frame}[t,fragile] \frametitle{Sentinel in code}
  \begin{columns}[T]
    \begin{column}{0.56\textwidth}
      \lstcode[basicstyle=\ttfamily\tiny]{codeoutput/SentinelDemo.java}{Java}
    \end{column}
    \begin{column}{0.40\textwidth}
      \lstterminal[basicstyle=\ttfamily\scriptsize]{codeoutput/SentinelDemo.txt}
      \small
      \hCode{head} is created once and never moves. \hCode{add} always
      inserts between \hCode{head} and \hCode{head.next}. No null check
      anywhere, no branch for the first node.
    \end{column}
  \end{columns}
\end{frame}
% ~~~~~~~~~~~~~~~~~~~~~~~~~~~~~~~~~~~~~~ A frm

% ~~~~~~~~~~~~~~~~~~~~~~~~~~~~~~~~~~~~~~ V
\hSection[Algos]{O(1)-Space Algorithms}
% ~~~~~~~~~~~~~~~~~~~~~~~~~~~~~~~~~~~~~~ A




% ~~~~~~~~~~~~~~~~~~~~~~~~~~~~~~~~~~~~~~ V frm
\begin{frame}[t,fragile] \frametitle{Morris traversal}
  \begin{columns}[T]
    \begin{column}{0.56\textwidth}
      \lstcode[basicstyle=\ttfamily\tiny]{codeoutput/MorrisTraversal.java}{Java}
    \end{column}
    \begin{column}{0.40\textwidth}
      \lstterminal[basicstyle=\ttfamily\scriptsize]{codeoutput/MorrisTraversal.txt}
      \small
      In-order traversal of a binary tree in \hRemark{$O(1)$ extra space}.
      The trick: temporarily overwrite a null \hCode{right} pointer
      with a thread back to the predecessor, then remove it.
    \end{column}
  \end{columns}
\end{frame}
% ~~~~~~~~~~~~~~~~~~~~~~~~~~~~~~~~~~~~~~ A frm




% ~~~~~~~~~~~~~~~~~~~~~~~~~~~~~~~~~~~~~~ V frm
\begin{frame}[t,fragile] \frametitle{When Morris pays}
  \begin{columns}[T]
    \begin{column}{0.5\textwidth}
      \begin{block}{Ordinary in-order}
        \small
        Recursion or an explicit stack: $O(h)$ extra space, where $h$
        is the height. On a balanced tree that is $O(\log n)$; on a
        skewed tree, $O(n)$.
      \end{block}
    \end{column}
    \begin{column}{0.46\textwidth}
      \begin{block}{Morris}
        \small
        $O(1)$ extra space, but each edge is visited twice
        (once to thread, once to remove). The time is still $O(n)$;
        only the constant doubles.
      \end{block}
    \end{column}
  \end{columns}
  \vspace{0.4em}
  \begin{itemize}\small
    \item Use Morris when memory is the constraint and you can afford
          the extra constant factor: embedded systems, huge trees,
          or a library that must not allocate.
    \item Otherwise use the iterative stack version, which is faster
          and easier to get right.
  \end{itemize}
\end{frame}
% ~~~~~~~~~~~~~~~~~~~~~~~~~~~~~~~~~~~~~~ A frm

% ~~~~~~~~~~~~~~~~~~~~~~~~~~~~~~~~~~~~~~ V
\hSection[API]{API Lessons}
% ~~~~~~~~~~~~~~~~~~~~~~~~~~~~~~~~~~~~~~ A




% ~~~~~~~~~~~~~~~~~~~~~~~~~~~~~~~~~~~~~~ V frm
\begin{frame}[t,fragile] \frametitle{\texttt{LinkedList} implements \texttt{Deque}}
\begin{lstlisting}[language=Java, basicstyle=\ttfamily\small, frame=single, xleftmargin=2em]
public class LinkedList<E>
    extends AbstractSequentialList<E>
    implements List<E>, Deque<E>,
               Cloneable, Serializable
\end{lstlisting}
  \begin{itemize}\small
    \item One class, four interfaces: \hCode{List}, \hCode{Deque},
          \hCode{Cloneable}, \hCode{Serializable}.
    \item That is why it can be used as a list, a stack
          (\hCode{push}/\hCode{pop}), a queue (\hCode{offer}/\hCode{poll}),
          or a double-ended queue (\hCode{addFirst}/\hCode{addLast}).
    \item \hRemark{One class, four vocabularies}: the same node machinery
          satisfies every interface.
  \end{itemize}
\end{frame}
% ~~~~~~~~~~~~~~~~~~~~~~~~~~~~~~~~~~~~~~ A frm




% ~~~~~~~~~~~~~~~~~~~~~~~~~~~~~~~~~~~~~~ V frm
\begin{frame}[t,fragile] \frametitle{\texttt{SequencedCollection} (Java 21)}
\begin{lstlisting}[language=Java, basicstyle=\ttfamily\small, frame=single, xleftmargin=2em]
interface SequencedCollection<E> extends Collection<E> {
    E    getFirst();
    E    getLast();
    void addFirst(E);
    void addLast(E);
    E    removeFirst();
    E    removeLast();
    SequencedCollection<E> reversed();
}
\end{lstlisting}
  \begin{itemize}\small
    \item Implemented by \hCode{LinkedList}, \hCode{ArrayList},
          \hCode{ArrayDeque}, \hCode{LinkedHashSet}, \hCode{TreeSet},
          and others.
    \item Before Java 21, \hCode{List.get(0)} and \hCode{Deque.getFirst()}
          had no common interface. Now they do.
    \item \hRemark{\texttt{reversed()} returns a live view}, not a copy.
  \end{itemize}
\end{frame}
% ~~~~~~~~~~~~~~~~~~~~~~~~~~~~~~~~~~~~~~ A frm




% ~~~~~~~~~~~~~~~~~~~~~~~~~~~~~~~~~~~~~~ V frm
\begin{frame}[t,fragile] \frametitle{\texttt{ConcurrentLinkedDeque}}
  \begin{columns}[T]
    \begin{column}{0.5\textwidth}
      \begin{block}{The shape}
        \small
        Same node structure, but \hCode{head} and \hCode{tail} are
        \hCode{AtomicReference}s. Every update is a CAS.
      \end{block}
    \end{column}
    \begin{column}{0.46\textwidth}
      \begin{block}{Properties}
        \small
        \begin{itemize}
          \item lock-free (Michael--Scott style)
          \item \emph{weakly consistent} iteration: no snapshot
          \item allows \hCode{null}? No.
          \item unbounded
        \end{itemize}
      \end{block}
    \end{column}
  \end{columns}
  \vspace{0.4em}
  \small
  \hRemark{Use it when many threads push and pop from both ends}.
  For one producer and one consumer, a bounded
  \hCode{ArrayBlockingQueue} is usually faster.
\end{frame}
% ~~~~~~~~~~~~~~~~~~~~~~~~~~~~~~~~~~~~~~ A frm

% ~~~~~~~~~~~~~~~~~~~~~~~~~~~~~~~~~~~~~~ V
\hSection[Wrap]{Wrap-up}
% ~~~~~~~~~~~~~~~~~~~~~~~~~~~~~~~~~~~~~~ A




% ~~~~~~~~~~~~~~~~~~~~~~~~~~~~~~~~~~~~~~ V frm
\begin{frame}[t,fragile] \frametitle{Big ideas}
  \begin{itemize}
    \item \hCode{LinkedList} is a doubly linked list with a
          \hCode{size} field, \hCode{first}, and \hCode{last}.
          That is the whole data structure.
    \item \hCode{ArrayList} is one \hCode{Object[]} plus a
          \hCode{size} counter. The two classes share only the interface.
    \item A \textbf{sentinel node} removes the empty-list special case
          from every add and remove.
    \item \textbf{Morris traversal} does an in-order walk of a tree in
          $O(1)$ extra space by threading \hCode{right} pointers
          through the tree and removing them as it goes.
    \item \hCode{LinkedList} implements \hCode{List} and \hCode{Deque}:
          one class, four vocabularies.
    \item \hCode{SequencedCollection} (Java 21) unified the first/last
          operations that \hCode{List}, \hCode{Deque} and \hCode{Set}
          had duplicated.
    \item \hCode{ConcurrentLinkedDeque} applies the same node structure
          with CAS on \hCode{head} and \hCode{tail} to make it lock-free.
  \end{itemize}
\end{frame}
% ~~~~~~~~~~~~~~~~~~~~~~~~~~~~~~~~~~~~~~ A frm




% ~~~~~~~~~~~~~~~~~~~~~~~~~~~~~~~~~~~~~~ V frm
\begin{frame}[t,fragile] \frametitle{Check yourself}
  \begin{enumerate}\small
    \item Why is \hCode{Node} declared \hCode{private static} inside
          \hCode{LinkedList}? What would break if it were not static?
    \item Why does \hCode{LinkedList} cache \hCode{size} instead of
          computing it by walking?
    \item What does a sentinel node save you, and what does it cost?
    \item Under what conditions does Morris traversal use less memory
          than the iterative stack version?
    \item Why is \hCode{SequencedCollection} a better name than
          adding the methods to \hCode{List}?
    \item Why is \hCode{ConcurrentLinkedDeque} \emph{weakly consistent}
          instead of fully consistent? Give an example of an iteration
          that would differ.
  \end{enumerate}
\end{frame}
% ~~~~~~~~~~~~~~~~~~~~~~~~~~~~~~~~~~~~~~ A frm




% ~~~~~~~~~~~~~~~~~~~~~~~~~~~~~~~~~~~~~~ V frm
\begin{frame}[t,fragile] \frametitle{Practice}
  \begin{enumerate}\small
    \item Run \hCode{javap -p java.util.LinkedList} on your JDK and
          compare the field list with the one on slide~5.
    \item Rewrite \hCode{SentinelDemo} without a sentinel. Count the
          null checks you had to add.
    \item Add a sentinel to \hCode{DoublyLinked<T>} from the sophomore
          deck. How many branches disappear?
    \item Trace \hCode{MorrisTraversal} on the tree
          \hCode{4 -> 2 -> 6} by hand. Which node's right pointer is
          temporarily pointed where?
    \item Replace the plain \hCode{AtomicReference} in a small
          lock-free queue with \hCode{AtomicStampedReference}. Does
          ABA still happen?
    \item Implement \hCode{reversed()} for \hCode{SinglyLinked<T>}
          without copying. What complexity does it have?
  \end{enumerate}
\end{frame}
% ~~~~~~~~~~~~~~~~~~~~~~~~~~~~~~~~~~~~~~ A frm




% ~~~~~~~~~~~~~~~~~~~~~~~~~~~~~~~~~~~~~~ V frm
\begin{frame}[t,fragile] \frametitle{References}
  \begin{enumerate}\small
    \item Robert Lafore, \emph{Data Structures \& Algorithms in Java},
          2nd ed., Sams, 2002.
    \item Joshua Bloch, \emph{Effective Java}, 3rd ed.,
          Addison-Wesley, 2018. (Items on \hCode{equals},
          \hCode{hashCode}, and defensive copying.)
    \item Doug Lea, \emph{Concurrent Programming in Java},
          2nd ed., Addison-Wesley, 1999.
    \item Joseph H. Morris, \emph{Traversing binary trees simply and
          cheaply}, Information Processing Letters, 9(5):197--200, 1979.
          (The original Morris traversal paper.)
    \item Oracle, \emph{Java SE 21 API Specification}:
          \hCode{java.util.LinkedList},
          \hCode{java.util.ArrayList},
          \hCode{java.util.concurrent.ConcurrentLinkedDeque},
          \hCode{java.util.SequencedCollection}.
    \item JEP 431: \emph{Sequenced Collections} (Java 21).
  \end{enumerate}
\end{frame}
% ~~~~~~~~~~~~~~~~~~~~~~~~~~~~~~~~~~~~~~ A frm




% ~~~~~~~~~~~~~~~~~~~~~~~~~~~~~~~~~~~~~~ V frm
\begin{frame}[t,fragile]
  \begin{center}
    \vspace{4cm}
    \hRemark{\Huge Questions?}
  \end{center}
\end{frame}
% ~~~~~~~~~~~~~~~~~~~~~~~~~~~~~~~~~~~~~~ A frm




\end{document}